\chapter{When the instrument is the fault}
\label{ch:instrument}

Sixteen times during this project a measurement was wrong and the system was
not. This chapter reports them as a body of evidence rather than as a defect
list, because the pattern across them is the most transferable outcome of the
work.

Most engineering write-ups treat measurement as transparent: the apparatus
reports, and the report is the finding. That assumption failed here often
enough to be worth quantifying. In every case the wrong reading was
\emph{plausible}. None announced itself. Several would have entered this
thesis as results had a second check not contradicted them, and two would have
removed findings that are now central to it.

\section{The rule}

\begin{quote}
\textbf{A surprising failure is evidence about the instrument until the
instrument has been cleared.}
\end{quote}

Four corollaries, each learned from a specific case below.

\begin{enumerate}
  \item Every analysis script gets a \textbf{known-answer test}: an input whose
    answer is known independently, checked before the script's output is
    believed.
  \item Prefer \textbf{real recorded data} whose properties were established
    separately.
  \item Use synthetic data only where the ground truth is \textbf{constructed}
    (arithmetic, geometry), never where it is \textbf{rendered} (images,
    physics). Two points \SI{300}{\milli\metre} apart really are that far
    apart; a rendered image is only as good as the renderer.
  \item A result that \textbf{contradicts an earlier measurement} is an
    instrument check, not a finding, until the two are reconciled.
\end{enumerate}

\section{The sixteen cases}

\Cref{tab:instr} lists all sixteen. Five predate the final development period
and eleven occurred within it, which is itself informative: the rate did not
fall as the project matured, because each new measurement brought a new
instrument.

\begin{table}[htbp]
  \centering
  \caption{Every occasion on which a measurement was wrong and the system was
  not. \emph{Caught by} is the thing that actually revealed it, not the thing
  that should have.}
  \label{tab:instr}
  \footnotesize
  \begin{tabular}{@{}p{32mm}p{45mm}p{45mm}@{}}
    \toprule
    Looked like & Actually was & Caught by \\
    \midrule
    Noise floor exactly zero on all 14 channels
      & The recorder sampled a \SIrange{14.7}{17}{\hertz} sensor at
        \SI{50}{\hertz}, so \SIrange{82}{88}{\percent} of adjacent rows were
        identical
      & The figure being \emph{too} clean; over distinct updates the real
        floor is \SIrange{0.2}{5.4}{\degree} \\
    Isotropy swinging by a factor of two between runs
      & The statistic, not the measurement: min-over-directions is maximally
        sensitive to one tail failure
      & Repeating to $N=10$, where the inter-quartile range is
        \SI{0.000}{\metre} \\
    Object detection at \SIrange{0}{4}{\percent}
      & The renderer: flat-shaded primitives are outside the detector's
        training distribution
      & The same model scoring \numrange{0.89}{0.91} on a real photograph \\
    A carry path failing 0 of 19 placements
      & The path template carried toward the wearer, the one direction leaving
        the feasible band
      & Contradiction with an earlier result showing 19 working grip pairs \\
    An {AprilTag} renderer detecting nothing
      & A mirrored tag: {OpenCV}'s camera has $+y$ down, reversing the winding
      & A vertical flip of the same crop decoding immediately \\
    A flat \SI{22}{\degree} pose error
      & The ground-truth rotation applied in the wrong frame; the two forms
        \emph{commute} about $x$ and $y$
      & Testing a general in-plane axis: \SI{21.79}{\degree} against
        \SI{1.06}{\degree} \\
    2--3 processes surviving a clean shutdown
      & The check grepped bare names and matched its own command line
      & Matching installed executable paths, which a command line cannot
        contain \\
    A terminal GUI launching correctly
      & The test built the one environment where the bug cannot occur, and
        checked only for captions
      & Running the real entry point and asserting on exit status \\
    \midrule
    The gripper model wrong by \SI{52}{\milli\metre}
      & The measurement took the gap between link \emph{origins}, which sit
        behind the pad faces
      & The vendor's published \SI{85}{\milli\metre} stroke: corrected, it
        lands at \SI{84.84}{\milli\metre} \\
    A wrist rotation of \SI{0}{\degree}, twice
      & First reading the anchor after the arm had moved onto the grasp pose;
        then reading it when a \emph{previous run} had left it there
      & Contradiction with an independently known \SI{169.7}{\degree} \\
    A known-answer check passing
      & It parsed zero rows, and \texttt{all()} over an empty collection is
        true
      & The output table being empty while the check said pass \\
    A health classifier missing an injected fault
      & The injected fault was a \SI{30}{\degree} \emph{circular} jump,
        correctly below the \SI{60}{\degree} threshold
      & Asking why only one of two injected faults was detected \\
    {RViz} failing to render inside a {Qt} container
      & The capture fired after the probe window had closed
      & Re-capturing during the window's lifetime \\
    A scene fingerprint reporting every object vanished
      & The sweep began before discovery matched the detector and received
        \emph{zero} messages
      & Running it twice on an unchanged scene \\
    A boundary probe reporting a wall that was not there
      & The test commanded an orientation unreachable at 0 of 6 points on that
        line
      & Instrumenting the node's own probe trace \\
    Nine sling clips with no object
      & Two faults at once: a stale \SI{350}{\milli\metre} length making the
        sling geometrically impossible, then floors calibrated at the wrong
        image scale
      & Looking at a frame, and the overlay reading \texttt{sag=0mm} \\
    \bottomrule
  \end{tabular}
\end{table}

\section{Four mechanisms}

Sorted by mechanism rather than by symptom, the sixteen fall into four groups.
None is specific to robotics.

\paragraph{Absence read as a value.} A known-answer check that passed on zero
parsed rows; a scene sweep that received no messages and concluded every
object had gone; a noise floor of zero produced by a recorder outrunning its
sensor. In each case code conflated \emph{no data} with \emph{a measurement of
zero}. This is the most dangerous group, because absence is silent by
definition and the resulting value is usually the most reassuring one
available.

\paragraph{The harness constructs the condition it tests.} A GUI test that
allocated the terminal in which the bug cannot occur; a detection rate taken
through a renderer the detector was never trained on; a fault injected at
\SI{30}{\degree} against a \SI{60}{\degree} threshold. A test written from the
same understanding as the code inherits that understanding's errors.

\paragraph{State left over from a previous run.} An anchor read while the arms
were still parked where the last run left them; a reachability sweep taken
with the arms \SIrange{155}{166}{\degree} from home; a boundary probe firing
from the previous run's end point. Scripts are written as though they begin
from nothing, and they do not.

\paragraph{Order of operations inside the measurement.} Reading a reference
after perturbing it, and screenshotting after the window closed. Both are easy
to do and invisible in the output afterwards.

\section{What actually caught them}

The defence that worked was not care. It was \textbf{redundancy}: having an
independently known quantity to check against.

\begin{table}[htbp]
  \centering
  \caption{What revealed each of the sixteen.}
  \label{tab:caught}
  \begin{tabular}{@{}lr@{}}
    \toprule
    Revealed by & Cases \\
    \midrule
    An independently known quantity contradicting the reading & \textbf{9} \\
    Inspection of the code or the method & 3 \\
    Looking at the raw artefact, usually an image & 4 \\
    \bottomrule
  \end{tabular}
\end{table}

Nine of sixteen were caught because something else already had an answer: a
vendor's published stroke, a published detection score, an earlier
incompatible result, an arithmetic prediction. Only three were caught by
inspection, and those were the shallowest.

This has a direct design consequence, and it is the recommendation this
chapter exists to make. \textbf{Build the cross-check in rather than trying
harder.} A measurement with no independent reference is not merely
unvalidated; it is structurally unable to be validated, and no amount of
diligence changes that. Where a project cannot obtain a second reference for a
quantity, the honest response is to mark it unverified rather than to raise
confidence in the single measurement.

\section{Consequences for this thesis}

Every quantitative claim in this document carries a maturity label and a
traceable measurement, and the audit records for each whether its instrument
was validated. Of thirty-four figures audited, fifteen are validated against an
independently known answer, fifteen are arithmetic or graph inspection with no
instrument to mis-calibrate, three are unverified and are not reported as
results, and one is structurally undefined in the available environment.

One capability was \textbf{cut rather than shipped} on this basis. A
boundary-warning node produced a correct \SI{88}{\milli\metre} lead in three of
five runs and a false alarm in the other two, and the cause survived
elimination of every candidate mechanism. A warning that is wrong teaches the
operator to ignore it, which spends the credibility on which every other alarm
in the system depends. Suppressing the symptom without knowing the cause was
available and was refused.

\todo{CONFIRM: whether this chapter belongs in the main body or as an
appendix. It is a methods contribution rather than a result, and its length is
justified only if the examiner reads it as such.}
